\begin{description}
\item[(T01.1)]
Since the two detectors make an angle with each other, photons from the
source will not be incident on the detectors at the same angle. Using this
the source can localized.
% sic: "can localized" as printed in the source.

Let $\widehat{n_1}$ be the unit vector perpendicular to the plane of
D$_1$ and $\widehat{n_2}$ be the unit vector perpendicular to the plane of
D$_2$. We have,
\begin{align*}
\widehat{n_1} &= \frac{1}{2}\widehat{x} + \frac{\sqrt{3}}{2}\widehat{y}
  && \text{(0.5 marks)} \\
\widehat{n_2} &= -\frac{1}{2}\widehat{x} + \frac{\sqrt{3}}{2}\widehat{y}
  && \text{(0.5 marks)}
\end{align*}
Let $\widehat{s} = s_x\widehat{x} + s_y\widehat{y}$ be the unit vector in
the direction of the source and $F$ be the flux of the source. The
effective area of capture for a detector will be
\[ A_{\mathrm{eff}} = A(|\widehat{s} \cdot \widehat{n}|) \]
\hfill(0.5 marks)
% FIGURE NEEDED: sketch of parallel rays hitting a tilted flat detector,
% showing unit vectors n-hat and s-hat and the projected effective area
% A_eff; 2025_theory_solution.pdf page 2.

Hence, the power received is,
\[ P = AF(|\widehat{s} \cdot \widehat{n}|) \]
\hfill(1.0 marks)\\
and therefore,
\begin{align*}
P_1 &= AF(|\widehat{s} \cdot \widehat{n_1}|) \\
P_2 &= AF(|\widehat{s} \cdot \widehat{n_2}|)
\end{align*}
Hence,
\begin{align*}
\widehat{s} \cdot \widehat{n_1}
  &= \frac{1}{2}s_x + \frac{\sqrt{3}}{2}s_y
   = \frac{2.70 \times 10^{-10}}{AF} \\
\widehat{s} \cdot \widehat{n_2}
  &= -\frac{1}{2}s_x + \frac{\sqrt{3}}{2}s_y
   = \frac{4.70 \times 10^{-10}}{AF}
\end{align*}
Solving,
\begin{align*}
s_x &= -\frac{2.00 \times 10^{-10}}{AF}
  && \text{(0.5 marks)} \\
s_y &= \frac{7.40 \times 10^{-10}}{\sqrt{3}AF}
  && \text{(0.5 marks)}
\end{align*}
Hence, the angle the direction vector makes with respect to the positive
$y$-axis is
\[ \eta = \tan^{-1}\left(-\frac{s_x}{s_y}\right) \]
\hfill(1.0 marks)

Then, $\eta = 25.1^\circ = 0.438$\,rad. \hfill(0.5 marks)

\item[(T01.2)]
The time delay occurs because the satellites S$_1$ and S$_2$ are located at
different distances from the source. Light takes longer time to reach the
satellite that is farther (here, say S$_1$). Let $\Delta t = t_1 - t_2$.

% FIGURE NEEDED: geometry of satellites S1 and S2 separated by 2r with the
% source direction making angle theta with the line joining them, path
% difference c*Delta-t marked; 2025_theory_solution.pdf page 3.

If the direction of the source makes an angle $\theta$ (see the figure
above) with the line joining telescopes S$_1$ and S$_2$ we have the path
difference ($\Delta d$) given as,
\[ \Delta d = c\Delta t \]
\hfill(0.5 marks)\\
From the figure above the path difference can be written as,
\[ \Delta d = 2r\cos\theta \]
\hfill(1.0 marks)\\
Therefore,
\[ \cos\theta = \frac{c\Delta t}{2r} \]
Since the measurement of $\Delta t$ has some uncertainty, that is
$\Delta t \in [9.9, 10.1]$\,ms, we get a range $[\theta_1, \theta_2]$ for the
possible values of $\theta$.
\begin{align*}
\theta_1 &= \cos^{-1}\left(\frac{c(10.1\,\mathrm{ms})}{2r}\right)
  = 77.51^\circ
  && \text{(0.5 marks)} \\
\theta_2 &= \cos^{-1}\left(\frac{c(9.9\,\mathrm{ms})}{2r}\right)
  = 77.76^\circ
  && \text{(0.5 marks)}
\end{align*}

% FIGURE NEEDED: unit sphere with two cones of half-angles theta_1 (blue)
% and theta_2 (red) about the baseline axis, bounding an annulus on the
% sphere where the source may lie; 2025_theory_solution.pdf page 3.

The source may lie anywhere in the annulus bounded by the red and blue
circles on the unit sphere above, whose area is given by
\[ 2\pi(\cos\theta_1 - \cos\theta_2) \]
\hfill(1.5 marks)\\
Hence the fraction of sky where the source might lie is
\[ f = \frac{1}{2}(\cos\theta_1 - \cos\theta_2) \]
\hfill(0.5 marks)
\[ f = 2.13 \times 10^{-3} \]
\hfill(0.5 marks)
\end{description}
